\documentclass[aps,prresearch,twocolumn,superscriptaddress,nofootinbib]{revtex4-2}
\usepackage{amsmath,amssymb,graphicx,booktabs}
\newcommand{\ii}{\mathrm{i}}
\newcommand{\tr}{\operatorname{tr}}

\begin{document}

\title{Trajectory-resolved retarded modular response of programmable topological domain walls}
\author{H1 numerical campaign}
\date{\today}

\begin{abstract}
We define a trajectory-resolved handedness diagnostic by applying an infinitesimal local
charge-phase kick to a reduced Gaussian state and evolving the resulting tangent in
modular time. The calculation compares hard support-terminated and soft explicit domain
walls at two controller strengths. A localized single-particle packet supplies a separate
visual corroboration. Modular time is generated by the prepared state's entanglement
Hamiltonian and is not monitored-circuit time or a quantized transport coefficient.
\end{abstract}

\maketitle

\section{Campaign and reduced state}

The production matrix fixes $N_x=20$, $N_y=40$, $n_{\rm shell}=1$,
$\alpha_2=30$, and $\alpha_1\in\{1,3\}$. We compare a hard wall with
\texttt{dw\_truncation=true} and \texttt{meas\_slab\_only=true} against a soft
wall with both flags false. Each case contains $S=25$ independent Haar-random
half-filled Slater trajectories, perfect correction, Born sampling, random serial order,
and 80 circuit cycles. Observations are made at
$t\in\{40,48,56,64,72,80\}$.

At every trajectory-checkpoint we trace out
$[0,N_x)\times[0,N_y/2)$ and retain
\begin{equation}
 A=[0,N_x)\times[N_y/2,N_y).
\end{equation}
With the one-body correlation convention
$C_{ij}=\langle\hat c_i^\dagger\hat c_j\rangle$, diagonalize
\begin{equation}
 C_A=U\,\operatorname{diag}(\nu_a)U^\dagger,
 \qquad
 h_A=\log\frac{\mathbf 1_A-C_A}{C_A}.
\end{equation}
Occupations are clipped at $\epsilon=10^{-8},10^{-10},10^{-12}$; the primary
choice is $10^{-10}$. One $h_A$ is used for all sources at a given
trajectory-checkpoint. Neither $C_A$ nor $h_A$ is averaged before constructing the
nonlinear observable.

\section{Retarded modular response}

For ten distinct random cells on each wall, freshly sampled at every
trajectory-checkpoint, define the two-orbital charge projector
\begin{equation}
 P_s=\sum_{\mu=A,B}\lvert x_w,y_s,\mu\rangle
 \langle x_w,y_s,\mu\rvert.
\end{equation}
The many-body kick $\exp(-\ii\lambda\hat N_s)$ produces the linear tangent
\begin{align}
 X_s(0)&=-\ii[P_s,C_A],\\
 X_s(\tau)&=e^{-\ii h_A\tau}X_s(0)e^{\ii h_A\tau},\\
 \chi_s^R(\boldsymbol r,\tau)&=\tr[P_{\boldsymbol r}X_s(\tau)].
\end{align}
This response is signed and number conserving:
$\sum_{\boldsymbol r}\chi_s^R(\boldsymbol r,\tau)=0$.
We evaluate $\tau=0,0.05,\ldots,4$.

For a wall window of periodic half-width $d$, define the source-relative dipole
\begin{equation}
 D_{w,s}(\tau)=\sum_{|x-x_w|_{\rm per}\le d}\sum_y
 (y-y_s)\chi_s^R(x,y,\tau).
\end{equation}
The primary $d=2$ estimator first aligns and averages the ten source responses within
one wall and trajectory. A linear fit on $\tau\in[0.1,0.8]$ gives $v_w$, followed by
\begin{equation}
 H=\frac{v_{x=5}-v_{x=15}}{2}.
\end{equation}
The six checkpoint values are averaged inside each trajectory. Confidence intervals
bootstrap only the 25 independent trajectories; sources and checkpoints are correlated
internal averages. Hard--soft and $\alpha_1$ contrasts use unpaired trajectory
bootstraps.

\section{Packet corroboration}

The supplemental packet begins in
\begin{equation}
 \lvert\phi_s(0)\rangle=\frac{1}{\sqrt2}
 \left(\lvert x_w,y_s,A\rangle+\lvert x_w,y_s,B\rangle\right)
\end{equation}
and evolves as $\lvert\phi_s(\tau)\rangle=e^{-\ii h_A\tau}\lvert\phi_s(0)\rangle$
for $\tau\in[0,8]$. Its positive probability, wall retention, displacement, and fitted
velocity test whether modular propagation remains wall localized. This artificial
one-particle probe is not a charge response and is not used as the primary H1 estimator.

\section{Finite-window Fourier diagnostic}

After projecting onto the primary wall window, the source-aligned response is transformed
according to
\begin{equation}
 \chi_w(k,\omega)=\Delta\tau\sum_{j,d_y}
 e^{\ii\omega\tau_j-\ii k d_y}
 W(\tau_j)e^{-\eta\tau_j}\chi_w(d_y,\tau_j),
\end{equation}
where $W$ is a Tukey taper and $\eta=0.35$. We repeat the extraction at
$\eta=0.20,0.50$. Ridge maxima on $0.2\le\omega\le2.5$ and
$0<|k|\le2.6$ are fitted to $\omega_{\rm peak}=b+v_F|k|$. Because the retained
interval has a finite open aperture, this is an exploratory consistency check rather
than a momentum-conserving susceptibility or controlled conformal dispersion.

\section{Results and diagnostics}

\begin{table}[b]
\caption{Late-time handed response. Each entry averages six checkpoints within each of
$S=25$ trajectories before bootstrap resampling.}
\centering
\IfFileExists{analysis_outputs/h1_results_table.tex}{%
\input{analysis_outputs/h1_results_table.tex}}{Production analysis pending.}
\end{table}

\begin{figure*}[t]
\centering
\IfFileExists{analysis_outputs/h1_real_space_retarded_response.pdf}{%
\includegraphics[width=0.95\textwidth]{analysis_outputs/h1_real_space_retarded_response.pdf}}{}
\caption{Source-aligned real-space retarded modular response. The production caption must
report the displayed modular time, $S=25$, trajectory bootstrap intervals, and the
source/checkpoint averaging order.}
\end{figure*}

\begin{figure}[b]
\centering
\IfFileExists{analysis_outputs/h1_handed_response_checkpoints.pdf}{%
\includegraphics[width=\columnwidth]{analysis_outputs/h1_handed_response_checkpoints.pdf}}{}
\caption{Checkpoint dependence of the primary handed response. Error bars are 95\%
trajectory-bootstrap intervals; ten sources per wall are averaged within each
trajectory-checkpoint.}
\end{figure}

Static modular susceptibility is intentionally omitted: the CPU pilot found that it
primarily measured local compressibility and did not supply a propagation direction.
No physical pass/fail threshold is imposed on H1. The campaign delivers the raw
trajectory-resolved response, localization and conservation diagnostics, Fourier
characterization, and packet corroboration.

\end{document}
